\documentclass[11pt]{article}
\usepackage{raeez-math-template}
\numberwithin{equation}{section}
\title{The Third Homology of a Marked Cusp Attachment}
\author{Raeez Lorgat}
\date{}
\hypersetup{pdfsubject={Integral third homology and meridian mapping tori}}
\begin{document}
\maketitle
\begin{abstract}
Four meridian mapping tori determine the third homology of a specified
two-filling manifold. Their connecting classes have one integral relation.
A closed cochain on the actual attachment identifies this relation with
a primitive core three-cycle. We compute the complete degree-three
attaching matrix, including the two fibre classes.
It is unimodular, and the resulting marked cusp attachment has zero
integral third homology.
\end{abstract}
\section{The real attachment and its chains}

All homology groups have coefficients in $\mathbb Z$.
Let $F=\mathbb R^4/\Lambda$ with ordered lattice
$\Lambda=\bigoplus_{j=1}^4\mathbb Ze_j$.
Matrices act on columns. Path multiplication is right to left.
The mapping torus of $T$ identifies $(a,1)$ with $(Ta,0)$.
Its positive meridian is $v\mapsto[0,v]$.
All mapping-torus fundamental classes use interval-first orientation.

Let $E\to\Sigma$ be the flat $F$-bundle over the sphere with three disks removed.
It has a zero section and marked meridians
$\gamma_0\gamma_1\gamma_\infty=1$.
Its monodromies are
\[
T_0=\begin{pmatrix}0&-1&-1&0\\1&-1&0&0\\0&0&1&0\\0&0&0&1\end{pmatrix},
\qquad
T_1=\begin{pmatrix}-1&1&-1&2\\-1&0&-1&1\\1&1&2&-1\\0&0&0&1\end{pmatrix},
\qquad T_\infty=(T_0T_1)^{-1}.
\]
For $X_i=\mathbb R^3/\mathbb Z^3$, put
\[
P_0=\begin{pmatrix}0&0&1\\1&0&0\\0&1&0\end{pmatrix},\qquad
P_1=\begin{pmatrix}0&-1&-1\\1&0&0\\0&0&1\end{pmatrix},\qquad
(d_0,d_1)=(3,4).
\]
Define $G_i=(X_i\times\mathbb R)/([y,\tau+1]=[P_i y,\tau])$,
and let $N_i=G_i\times D^2$.
The ordered coordinate changes are
\begin{equation}\label{th:eq:B}
B_0=\begin{pmatrix}0&-1&-1&4\\0&0&-1&2\\1&1&1&-6\\0&0&0&1\end{pmatrix},
\qquad
B_1=\begin{pmatrix}1&1&0&-1\\0&1&0&-1\\0&-1&1&3\\0&0&0&1\end{pmatrix}.
\end{equation}
Write $(y_i,t)^t=B_i^{-1}a$, so $t=a_4$.
The boundary maps and their core restrictions are
\begin{equation}\label{th:eq:maps}
\begin{aligned}
j_0[a,v]&=([y_0,3t+v],e^{-2\pi it}),&
f_0(a)&=[y_0,3t],\\
j_1[a,v]&=([y_1,4t-v],e^{2\pi it}),&
f_1(a)&=[y_1,4t].
\end{aligned}
\end{equation}
Indeed $B_0^{-1}T_0B_0=\operatorname{diag}(P_0,1)$ and
$B_1^{-1}T_1B_1=\operatorname{diag}(P_1^{-1},1)$.
These identities verify the meridian seams.
The two circle matrices have determinant one, and $B_i$ is unimodular.
Thus $j_i$ are diffeomorphisms and $f_i$ are coverings of degrees $d_i$.
Use the induced boundary orientations to form
\[
W=E\cup_{j_0}N_0\cup_{j_1}N_1,\qquad \partial W=B_{T_\infty}.
\]
This is a specified compact real six-manifold.

The base retracts onto its two finite boundary circles joined by an arc.
Flat transport lifts this retraction to $E$.
Attach the two fillings using boundary collars and then contract their disk factors.
This gives a homotopy equivalence with the double mapping cylinder
\[
D=G_0\cup_{f_0}(F\times[0,1])\cup_{f_1}G_1.
\]
It retains the meridian homotopies $[y_0,3t+v]$ and $[y_1,4t-v]$.
In particular, the positive cusp path first traverses $\gamma_0^{-1}$
and then $\gamma_1^{-1}$, with the joining arc between them.

These maps admit finite compatible triangulations.
Intersect the finitely many $P_i$-translates of a rational torus subdivision.
Triangulate that intersection and use its prisms for $G_i$.
Cut the source of $f_i$ at integral values of $d_it$ and at inverse images of target faces.
Common refinements give cellular maps and the double mapping cylinder.
Every additional rational subtorus below can be included by the same refinement.
There is no requirement for one finite triangulation invariant under both $T_0$ and $T_1$.

We use the chain cone
\begin{equation}\label{th:eq:cone}
\begin{aligned}
\mathcal D_n&=C_n(G_0)\oplus C_n(G_1)\oplus C_{n-1}(F),\\
d(a_0,a_1,b)&=(da_0+f_{0*}b,\,da_1-f_{1*}b,\,-db).
\end{aligned}
\end{equation}
Here $C_*$ can be the compatible cellular chains or piecewise smooth singular chains.
Subdivision and prism comparisons identify the two diagrams.
The shifted term uses the negative of the interval-first cylinder orientation.
The resulting exact sequence contains
\begin{equation}\label{th:eq:sequence}
0\longrightarrow\operatorname{coker}J_3
\longrightarrow H_3(W)\xrightarrow{\partial}\ker J_2
\longrightarrow0,\qquad J_n=(f_{0*},-f_{1*}).
\end{equation}
The connecting map sends a cone cycle to $[b]$.

\section{Integral three-cycles in the finite cores}

Write $x_1,x_2,x_3$ for the separate ordered coordinate bases of $X_i$.
Let $A_i\in H_3(G_i)$ be the oriented fibre three-torus.
The invariant oriented planes have primitive bivectors
\[
w_0=x_1x_2-x_1x_3+x_2x_3,\qquad w_1=x_1x_2.
\]
Juxtaposition denotes exterior product.
Their primitive sublattices are
$\langle x_1-x_3,x_2-x_3\rangle$ and $\langle x_1,x_2\rangle$.
The matrices $P_i$ preserve these sublattices with determinant one.
Let $S_i$ be the fundamental class of the positive core mapping torus of that plane.

The core Wang sequence is
\[
0\longrightarrow
\operatorname{coker}(\bigwedge^3P_i-I)
\longrightarrow H_3(G_i)
\longrightarrow\ker(\bigwedge^2P_i-I)
\longrightarrow0.
\]
It follows directly from the interval-prism differential.
Since $\det P_i=1$, its left group is $\mathbb Z A_i$.
Solving the exterior-square equations gives $\mathbb Zw_i$ on the right.
The class $S_i$ maps to $w_i$.
Thus $(A_i,S_i)$ is an integral basis of $H_3(G_i)$.

In the adapted source basis $(x_1,x_2,x_3,t)$, the covering matrices are
\[
F^{\rm ad}_{0,3}=
\begin{pmatrix}1&0&0&0\\0&1&-1&1\end{pmatrix},\qquad
F^{\rm ad}_{1,3}=
\begin{pmatrix}1&0&0&0\\0&4&2&-2\end{pmatrix}.
\]
Their columns are the lexicographic exterior triples.
To verify the second rows, divide the covering interval into $d_i$ pieces.
On a mixed three-torus, the pieces contribute the norm
$\sum_{r=0}^{d_i-1}\bigwedge^2P_i^r$.
Moving the last interval to the first position contributes sign $+1$.
The two norm matrices are the columns $w_i$ times the rows
$(1,-1,1)$ and $(4,2,-2)$, respectively.
The first row on a mixed three-torus is zero:
the invariant volume form $dy_1\wedge dy_2\wedge dy_3$
has only two independent fibre derivatives there.
It integrates to one on $A_i$ and zero on $S_i$.
This excludes an additional fibre-volume coefficient.

Transforming by $\bigwedge^3B_i^{-1}$ gives the actual fibre maps
\begin{equation}\label{th:eq:F3}
F_{0,3}=\begin{pmatrix}1&6&2&-4\\0&3&1&-2\end{pmatrix},\qquad
F_{1,3}=\begin{pmatrix}1&-3&-1&1\\0&6&2&-4\end{pmatrix}.
\end{equation}
Consequently
\begin{equation}\label{th:eq:J3}
J_3=\begin{pmatrix}
1&6&2&-4\\0&3&1&-2\\-1&3&1&-1\\0&-6&-2&4
\end{pmatrix}
\end{equation}
in target basis $(A_0,S_0,A_1,S_1)$.
Its second column is three times its third.
Its first, third, and fourth columns, followed by $(0,0,0,1)^t$,
form a determinant-one matrix.
Hence
\begin{equation}\label{th:eq:c}
\operatorname{coker}J_3=\mathbb Zc,\qquad
c=[S_1],\qquad q=(0,2,0,1),\qquad q(c)=1.
\end{equation}
The primitive row $q$ annihilates $J_3$.

We also need the degree-two fibre relations.
The degree-two core bases consist of one fibre torus and the mapping torus
of $U_0=(1,1,1)^t$ or $U_1=(-1,-1,2)^t$.
The fibre quotient rows are $(1,-1,1)$ and $(2,1,-1)$.
The mixed-torus coefficients are the negative norms of $P_i$ on $H_1(X_i)$.
Invariant vertical two-forms exclude a residual fibre coefficient, as above.
This gives
\begin{equation}\label{th:eq:J2}
J_2=
\begin{pmatrix}
3&1&0&-2&0&0\\0&0&0&0&0&-1\\
-3&-1&0&2&-3&-1\\0&0&0&0&2&2
\end{pmatrix}.
\end{equation}
Its source basis is $(e_1e_2,e_1e_3,e_1e_4,e_2e_3,e_2e_4,e_3e_4)$.
An integral vector $y$ belongs to its kernel exactly when
\begin{equation}\label{th:eq:kerJ2}
y_5=y_6=0,\qquad 3y_1+y_2-2y_4=0.
\end{equation}

\section{Four actual meridian three-manifolds}

Use the ordered fibre basis
\[
\alpha_1=e_1-e_3,\quad \alpha_2=e_3-e_2,\quad
\beta_1=e_3,\quad\beta_2=e_4.
\]
It is unimodular, and
$T_\infty\alpha_j=\alpha_j$, $T_\infty\beta_j=\beta_j+\alpha_j$.
Consider the four primitive invariant subtori with oriented bivectors
\begin{equation}\label{th:eq:v}
v_1=\alpha_1\alpha_2,\quad
v_2=\alpha_1\beta_1,\quad
v_3=\alpha_2\beta_2,\quad
v_4=(\alpha_1+\alpha_2)(\beta_1+\beta_2).
\end{equation}
Let $V_j$ be the corresponding oriented torus and $M_j=[M(V_j)]$
its positive mapping-torus fundamental class in $H_3(B_{T_\infty})$.
The first restricted monodromy is the identity.
Each of the other three has matrix
$\left(\begin{smallmatrix}1&1\\0&1\end{smallmatrix}\right)$
in the displayed ordered basis.
Thus these are actual closed oriented three-manifolds.

The following chain construction retains their nontrivial seams.
Choose a finite oriented fundamental surface chain $u_j$ in $V_j$.
There is a piecewise smooth singular three-chain $a_j$, supported in $V_j$, with
\begin{equation}\label{th:eq:seam}
da_j=T_{\infty*}u_j-u_j.
\end{equation}
To construct it, take a common refinement of the triangulations for $u_j$
and $T_{\infty*}u_j$.
Let $p_j,p'_j$ be their subdivision prisms to its same oriented fundamental chain $u'_j$.
Then $dp_j=u'_j-u_j$ and $dp'_j=u'_j-T_{\infty*}u_j$.
Set $a_j=p_j-p'_j$.
The prism maps stay in $V_j$.
This construction uses surface-chain subdivision, without a homotopy between
the identity map and the restricted unipotent automorphism.

Put $u_j^-=T_{0*}^{-1}u_j$.
Let $R_i$ be the first three rows of $B_i^{-1}$.
The two inverse-meridian prisms, extended linearly over surface chains, are
\begin{equation}\label{th:eq:prisms}
Q_0(v,a)=[R_0a,3a_4-v],\qquad
Q_1(v,a)=[R_1a,4a_4+v].
\end{equation}
They define actual smooth maps $[0,1]\times F\to G_i$.
The period relations make these formulas independent of the lift of $a$.
Their endpoint identities give
\[
dQ_0(u_j)=f_{0*}(u_j^--u_j),\qquad
dQ_1(u_j^-)=f_{1*}(T_{\infty*}u_j-u_j^-).
\]
Thus the actual global cycle is
\begin{equation}\label{th:eq:globalcycle}
Z_j=\bigl(Q_0(u_j),\,Q_1(u_j^-)-f_{1*}a_j,\,u_j-u_j^-\bigr)
\in\mathcal D_3.
\end{equation}
Substituting \eqref{th:eq:seam} into \eqref{th:eq:cone} gives $dZ_j=0$
in each component, with the displayed signs.
Cutting $M(V_j)$ along its marked fibre gives its interval prism and the seam correction $-a_j$.
The spine retraction takes these pieces to \eqref{th:eq:globalcycle}.
Therefore $[Z_j]=b_{\infty*}M_j$ for the actual boundary inclusion.

The connecting vectors are
$\partial[Z_j]=(I-\bigwedge^2T_0^{-1})v_j$.
Their complete matrix is
\begin{equation}\label{th:eq:D}
D_2=
\begin{pmatrix}
-1&0&0&-1\\3&2&0&3\\0&0&2&3\\
0&1&0&0\\0&0&0&0\\0&0&0&0
\end{pmatrix}.
\end{equation}
Its columns $d_j$ satisfy \eqref{th:eq:kerJ2}.
The three vectors $d_1,d_2,d_3-d_4$ form an integral basis of that kernel.
Indeed, in those coordinates an element is
\[
(-r+t,\,3r+2s-3t,\,-t,\,s,\,0,\,0),
\qquad (r,s,t)\in\mathbb Z^3,
\]
and \eqref{th:eq:kerJ2} gives the unique inverse coordinates.
Solving $D_2n=0$ over $\mathbb Z$ gives
\begin{equation}\label{th:eq:r}
\ker D_2=\mathbb Z(-2,0,-3,2)^t.
\end{equation}
The corresponding combination of the $Z_j$ lies in the core subgroup $\mathbb Zc$.
Its coefficient is the remaining integral question.

\section{A closed cochain determines the primitive relation}

On $X_0$ and $X_1$, respectively, put
\[
\omega_0=dy_1\wedge dy_2-dy_1\wedge dy_3+dy_2\wedge dy_3,\qquad
\omega_1=2dy_1\wedge dy_2+dy_1\wedge dy_3-dy_2\wedge dy_3.
\]
Each is $P_i$-invariant.
Define the closed core forms
\begin{equation}\label{th:eq:theta}
\theta_0=\tfrac23\,d\tau\wedge\omega_0,\qquad
\theta_1=\tfrac12\,d\tau\wedge\omega_1.
\end{equation}
The integrals of $\omega_0,\omega_1$ on the invariant planes $w_0,w_1$
are $3$ and $2$.
Consequently the pair of forms evaluates on
$(A_0,S_0,A_1,S_1)$ as $(0,2,0,1)$.
Its restriction to the core quotient is the primitive homomorphism $q$.

The actual pullbacks agree as differential forms on $F$:
\begin{equation}\label{th:eq:eta}
f_0^*\theta_0=f_1^*\theta_1=\eta,\qquad
\eta=6\,da_1\wedge da_2\wedge da_4
 +2\,da_1\wedge da_3\wedge da_4
 -4\,da_2\wedge da_3\wedge da_4.
\end{equation}
Here $a_1,\ldots,a_4$ are the real coordinates in the $e$ basis.
This equality follows by substituting \eqref{th:eq:B} into \eqref{th:eq:theta}.
It is an equality of constant forms, including all coefficients.

On the piecewise smooth chain cone define
\[
\mathcal L(a_0,a_1,b)=\int_{a_0}\theta_0+\int_{a_1}\theta_1.
\]
For a boundary, Stokes' formula gives zero on the two core differentials.
The remaining terms are
$\int_b f_0^*\theta_0-\int_b f_1^*\theta_1=0$.
Thus $\mathcal L$ is a real homology functional.
It need not be integral on every global class.
Its restriction to the integral core subgroup is exactly $q$.

Each $a_j$ in \eqref{th:eq:seam} is supported in a two-dimensional torus.
Therefore $\int_{a_j}\eta=0$.
In the first prism, the meridian derivative of the core parameter is $-1$.
In the second prism it is $+1$.
Their fibre derivatives give
\begin{equation}\label{th:eq:evaluation}
\mathcal L(Z_j)
=-\tfrac23\,\omega_0\bigl((\bigwedge^2R_0)v_j\bigr)
 {}+\tfrac12\,\omega_1\bigl((\bigwedge^2(R_1T_0^{-1}))v_j\bigr).
\end{equation}
In the $e$ exterior basis, this evaluation row is
\[
\lambda=(-\tfrac12,-\tfrac16,-\tfrac32,\tfrac13,0,-1).
\]
Its four values on \eqref{th:eq:v} are
\begin{equation}\label{th:eq:values}
(\mathcal L(Z_1),\mathcal L(Z_2),\mathcal L(Z_3),\mathcal L(Z_4))
=(0,-\tfrac16,-1,-2).
\end{equation}
The fractional second value is compatible with the stated real coefficient field.
Using \eqref{th:eq:r}, the integral core class
$-2[Z_1]-3[Z_3]+2[Z_4]$ has evaluation $-1$.
Since $q(c)=1$ on the free core subgroup, this proves the integral relation
\begin{equation}\label{th:eq:relation}
-2[Z_1]-3[Z_3]+2[Z_4]=-c
\quad\hbox{in }H_3(W;\mathbb Z).
\end{equation}
In particular, the relation produces a primitive core class.
No intersection or linking pairing is needed for this calculation.

\section{The integral attaching matrix}

Set
$t_1=[Z_1]$, $t_2=[Z_2]$, and $t_3=[Z_3]-[Z_4]$.
Their connecting classes form the integral basis of $\ker J_2$ given above.
They therefore split \eqref{th:eq:sequence}.
Together with $c$, they are an integral basis of $H_3(W)$.
Equation~\eqref{th:eq:relation} now gives
\begin{equation}\label{th:eq:meridianmatrix}
b_{\infty*}(M_1,M_2,M_3,M_4)=
\begin{pmatrix}
1&0&-2&-2\\0&1&0&0\\0&0&-2&-3\\0&0&1&1
\end{pmatrix}
\end{equation}
in the target basis $(t_1,t_2,t_3,c)$.
The determinant is one.
This proves both the image and kernel over the integers.

We state the marked cusp input for the complete attachment.
Let $N_c$ be a compact smooth cusp neighborhood with boundary $B=B_{T_\infty}$.
Let $k:B\to N_c$ be its actual marked inclusion.
Assume its degree-three homology is free of rank two, and that
\begin{equation}\label{th:eq:local}
H_3(B)=
\bigoplus_{j=1}^4\mathbb ZM_j\oplus\mathbb ZC_1\oplus\mathbb ZC_2,
\qquad
k_*M_j=0,\quad k_*C_i=n_i,
\end{equation}
where $(n_1,n_2)$ is an integral basis of $H_3(N_c)$.
The fibre classes $C_1,C_2$ are the actual oriented three-tori
$\alpha_1\beta_1\beta_2$ and $\alpha_2\beta_1\beta_2$.
These are the local data of the marked toric specialization with periods
$(qc_1,c_2),(c_3,qc_4)$ and its retained section.
The local identification must preserve the actual meridian homotopy.

For clarity, the boundary-group splitting in \eqref{th:eq:local}
has a direct integral check.
The four $v_j$ form a basis of $\ker(\bigwedge^2T_\infty-I)$.
The fibre quotient in degree three is freely generated by $C_1,C_2$.
The actual $M_j$ project to $v_j$ in the Wang sequence.
For the marked toric collapse, their images have two-dimensional carriers,
and the fibre map sends $C_i$ to $n_i$.
These geometric properties are the local input in \eqref{th:eq:local}.

The fibre inclusion into $W$ lands in the core quotient.
The evaluation \eqref{th:eq:eta} gives
$b_{\infty*}C_1=2c$ and $b_{\infty*}C_2=4c$.
Consequently the full degree-three attaching map is
\begin{equation}\label{th:eq:full}
(b_{\infty*},-k_*)_3=
\begin{pmatrix}
1&0&-2&-2&0&0\\
0&1&0&0&0&0\\
0&0&-2&-3&0&0\\
0&0&1&1&2&4\\
0&0&0&0&-1&0\\
0&0&0&0&0&-1
\end{pmatrix}.
\end{equation}
The source basis is $(M_1,M_2,M_3,M_4,C_1,C_2)$.
The target basis is $(t_1,t_2,t_3,c,n_1,n_2)$.

\begin{theorem}\label{th:thm:third}
For the specified real fillings,
$H_3(W;\mathbb Z)=\mathbb Z^4$, with basis $(t_1,t_2,t_3,c)$.
The actual meridian map \eqref{th:eq:meridianmatrix} is an integral isomorphism.
Under the marked cusp data \eqref{th:eq:local}, the full map
\eqref{th:eq:full} is an integral isomorphism.
Both its kernel and cokernel vanish.
\end{theorem}
\begin{proof}
The global sequence, the basis of its quotient, and the primitive relation
give the stated integral basis.
The meridian determinant is one.
The last two diagonal entries of the full block matrix are both $-1$,
so its determinant is also one.
\end{proof}

To determine $H_3(Y)$ for $Y=W\cup_BN_c$, the preceding attaching map
in degree two must also be injective.
Here is the exact local input and its integral verification.
Assume $k_{*,2}$ is onto with kernel the two angular meridian tori
$z_i=M(\alpha_i)$.
The one-dimensional version of the actual prisms \eqref{th:eq:prisms}
gives their connecting images $(I-T_0^{-1})\alpha_i=e_1,2e_1$.
The difference $z_2-2z_1$ has fibre circle
$\delta=-2e_1-e_2+3e_3$, fixed by both finite monodromies.
Its core images are the negative $U_0$ mapping torus
and the positive $U_1$ mapping torus.
In the degree-two core basis, their sum is $c_2=(0,-1,0,1)^t$.
Columns $2,5,6$ of \eqref{th:eq:J2}, followed by $c_2$, have determinant one.
The degree-one fibre map, obtained directly from \eqref{th:eq:maps}, is
\[
J_1=\begin{pmatrix}0&0&1&6\\0&0&0&3\\0&-1&-1&2\\0&0&0&-4\end{pmatrix}.
\]
Its rows give the integral kernel $\mathbb Ze_1$.
Thus the two actual angular meridian tori map isomorphically onto $H_2(W)$.
If an element of $H_2(B)$ is killed by both $b_{\infty*}$ and $k_*$,
it belongs to their two-dimensional meridian subgroup and is zero.
This proves the needed injectivity of $(b_{\infty*},-k_*)_2$.

\begin{corollary}\label{th:cor:Y}
For the specified marked toric attachment, with the local data just stated,
\[
H_3(Y;\mathbb Z)=0.
\]
In particular, this group has no torsion.
\end{corollary}
\begin{proof}
The Mayer--Vietoris sequence gives
\[
\operatorname{coker}(b_{\infty*},-k_*)_3
\longrightarrow H_3(Y)
\longrightarrow\ker(b_{\infty*},-k_*)_2.
\]
Both outer groups vanish by the actual matrices and the preceding argument.
\end{proof}

The complex-family question retains a separate geometric comparison.
It concerns $Y(u,a_0,a_1)$ for all sufficiently small nonzero complex $u$
and primitive increments
$a_0\in\frac13\Lambda^{T_0}$, $a_1\in\frac14\Lambda^{T_1}$,
with $|12e_4^*(a_0+a_1)|=1$.
The desired conclusion is a diffeomorphism of the underlying smooth manifold with $S^6$.
The present finite maps use $a_0=e_4/3$ and $a_1=-(e_1+e_4)/4$.
Application requires an actual comparison of the punctured complex family with $E$,
relative to all three marked boundaries and the section.
More precisely, if $\mathcal E_u\to\Sigma$ is that restricted family,
the next map is a fibre-preserving diffeomorphism
$\Phi_u:\mathcal E_u\to E$ carrying its section to the zero section.
Its three boundary restrictions must equal the specified period markings.
For each local filling, an extension
$\Psi_i:\mathcal N_{u,i}\to N_i$ must satisfy
$\Psi_i j_i^{\rm geom}=j_i\Phi_u$ on the corresponding boundary.
The analogous equation must hold for the marked cusp filling.
Its finite lifted translations must extend smoothly and equivariantly,
with proper quotients, the specified periods, and conormal comparisons
$e=s^{-2}e'$ and $e=s^{-1}e'$.
Those maps, the other permitted affine increments, and the remaining topology
are separate from the third-homology calculation.
The toric and complex-geometric constructions motivating this question are
described in \cite[Sections 2--4]{THengel}.

\begin{thebibliography}{9}
\bibitem{THengel}
P.~Engel, \emph{Complex structures on $S^6$},
\href{https://philip-engel.github.io/S6.pdf}{author manuscript},
Propositions 2.7--2.9 and Sections 3--4.
\end{thebibliography}

\end{document}
